% Part: computability
% Chapter: computability-theory
% Section: russells-paradox

\documentclass[../../../include/open-logic-section]{subfiles}

\begin{document}

\olfileid[tr]{cmp}{thy}{rus}
\olsection{Russell Paradoksuyla Karşılaştırma}

Bu bölümdeki savları Russell paradoksuyla karşılaştırmak öğreticidir:
\begin{enumerate}
\item Russell paradoksu: $S=\Setabs{x}{x\notin x}$ olsun. O zaman
  $S\in S$ ancak ve ancak $S\notin S$ olur; bu bir çelişkidir.

  \emph{Sonuç:} Böyle bir $S$ kümesi yoktur. ``Bütün kümelerin kümesi''nin
  varlığını varsaymak küme kuramının öteki aksiyomlarıyla tutarsızdır.

\item Russell paradoksunun değiştirilmiş bir sürümü: $F$, bütün fonksiyonlar
  kümesinden $\{0,1\}$'e giden ve
  \[
    F(f)=\begin{cases}
      1 & \text{$f$, $f$'nin tanım kümesindeyse ve $f(f)=0$ ise}\\
      0 & \text{aksi durumda}
    \end{cases}
  \]
  ile verilen bir ``fonksiyon'' olsun. Benzer bir sav
  $F(F)=0$ ancak ve ancak $F(F)=1$ olduğunu gösterir; bu bir çelişkidir.

  \emph{Sonuç:} $F$ bir fonksiyon değildir. ``Bütün fonksiyonlar kümesi'',
  bir fonksiyonun tanım kümesi olamayacak kadar büyüktür.

\item Köşegenleştirme savı: $f_0,f_1,\dots$, kısmi hesaplanabilir
  fonksiyonların numaralandırması olsun ve $G\colon\Nat\to\{0,1\}$
  \[
    G(x)=\begin{cases}
      1 & \text{$f_x(x)\downarrow=0$ ise}\\
      0 & \text{aksi durumda}
    \end{cases}
  \]
  ile tanımlansın. $G$ hesaplanabilir olsaydı bazı $k$ için $f_k$ fonksiyonu
  olurdu. Fakat o zaman $G(k)=1$ ancak ve ancak $G(k)=0$ olurdu; bu bir
  çelişkidir.

  \emph{Sonuç:} $G$ hesaplanabilir değildir. Küme kuramı aksiyomlarına göre
  $G$'nin yine de bir fonksiyon olduğuna dikkat edin. Burada paradoks değil,
  yalnızca bir ayrımın açıklığa kavuşturulması vardır.
\end{enumerate}

Kısmi fonksiyonlar, hesaplanabilir fonksiyonlar, kısmi hesaplanabilir
fonksiyonlar vb. hakkında konuşmak kafa karıştırıcı olabilir. $\Nat$'den
$\Nat$'ye giden bütün kısmi fonksiyonlar kümesi büyük bir nesne
topluluğudur. Bunların bazıları tümel, bazıları hesaplanabilir, bazıları hem
tümel hem hesaplanabilir, bazıları da ikisi birden değildir. Varsayılan olarak
``fonksiyon'' dediğimizde tümel fonksiyonu kastettiğimizi unutmayın. Böylece
şu sınıflara sahibiz:
\begin{enumerate}
\item hesaplanabilir fonksiyonlar,
\item tümel olmayan kısmi hesaplanabilir fonksiyonlar,
\item hesaplanabilir olmayan fonksiyonlar,
\item ne tümel ne hesaplanabilir olan kısmi fonksiyonlar.
\end{enumerate}
Bu ayrımı düzenlemek için $\Nat$'den $\Nat$'ye giden bütün kısmi
fonksiyonları temsil eden büyük bir kare çizebilir ve içinde sırasıyla tümel
fonksiyonlara ve hesaplanabilir kısmi fonksiyonlara karşılık gelen, birbiriyle
kesişen iki bölgeyi işaretleyebilirsiniz. Ortaya çıkan bölgelerin her biri
için bir nesne betimleyip betimleyemeyeceğinizi görmek iyi bir alıştırmadır.

\end{document}

